\documentclass{homework}
% \usepackage{graphicx} % Required for inserting images
\usepackage{amsmath}
\usepackage{gensymb}
\usepackage{subfig}
\usepackage{float}
\usepackage{amsthm}


\class{ECEN 5448}
\title{Midterm 1}
\author{Sean Jennings}
\date{October 16}

\begin{document}
\maketitle

\section{Problem 1}

\begin{letterenum}
\item We define state variable $x = [x_1, x_2]^T\in \R^2$ such that $x_1 = \theta$
and $x_2 = \dot{\theta}$. Then, the system is nonlinear and can be represented by
the equations
\begin{align*}
    \dot{x} = \mat{\dot{x_1} \\ \dot{x_2}} = f(x, u) \\
    y = g(x, u)
\end{align*}
where $f: \R^2 \times \R \to \R^2$ is
\begin{equation*}
    f([x_1, x_2]^T, u) = \mat{
        x_2 \\
        -\frac{b}{I} x_2 - \frac{mg\ell}{I} \cos x_1
    } + \mat{0 \\ 1} u
\end{equation*}
and $g: \R^2 \times \R \to \R$ is
\begin{equation*}
    g([x_1, x_2]^T, u) = \ell \sin(x_1)
\end{equation*}

\item Since we assume that $\tau(t) = 0$ for all $t \geq 0$, we also have
    $u(t) = \tau(t) = 0$. Let $x^* = [\theta^*, \dot{\theta}^*]^T = [-\pi/2, 0]^T$.
    and $x^{**} = [\theta^{**}, \dot{\theta}^{**}]^T = [\pi/2, 0]^T$. To show
    that $x^*$ and $x^{**}$ are each equilibrium points with $u(t) = 0$,
    it suffices to show that $f(x^*, 0) = f(x^{**}, 0) = 0_{2\times 1}$, where
    $0_{2\times 1}$ is the zero vector in $\R^2$. Since
    $\cos(\pi/2) = \cos({-\pi/2}) = 0$, we have
    \begin{align*}
        f(x^*, 0) = f([-\pi/2, 0]^T, 0) = \mat{
            0 \\
            -\frac{b}{I} 0 - \frac{mg\ell}{I} \cos(-\pi/2)
        } + \mat{0 \\ 1} 0
        &= \mat{0 \\ 0}
    \end{align*}
    and similarly,
    \begin{align*}
        f_(x^{**}, 0) = f([\pi/2, 0]^T, 0) = \mat{0 \\ 0}.
    \end{align*}

    To compute the linearizations about these equilibrium points, we define
    two changes of variables $\delta x^* = x - x^*$
    and $\delta x^{**} = x - x^{**}$.
    
    We will also define $u^* = 0$ and $\delta u^* = u - u^* = u$ for completeness.
    But since $u(t) = 0$, we also have $\delta u^*(t) = 0$ for all $t$.

    Since $x^*$ and $x^{**}$ are constant, we have
    $\frac{d}{dt}(\delta x^*) = \frac{d}{dt}(\delta x^{**}) = f(x, u)$. Starting
    with $\delta x^*$, we take the Taylor expansion of $f$ about the equilibrium 
    point $\delta x^* = 0$:
    \begin{align*}
        \frac{d}{dt}(\delta x^*) &= f(x^*, u^*) +
            \frac{\delta f}{\delta x}(x^*, u^*) \delta x^*
            + \frac{\delta f}{\delta u}(x^*, u^*) \delta u^*
            + \mathcal{O}((\delta x^*)^2) + \mathcal{O}((\delta u^*)^2) \\
            &= \frac{\delta f}{\delta x}(x^*, u^*) \delta x^* + \mathcal{O}((\delta x^*)^2)
    \end{align*}
    where the second equality holds because $f(x^*, u^*) = 0$ and
    $\delta u^* = 0$. Similarly, for the second equilibrium point we have
    \begin{equation*}
        \frac{d}{dt}{(\delta x^{**})} = \frac{\delta f}{\delta x} (x^{**}, u^*)\delta x^{**}
            + \mathcal{O}((\delta x^{**})^2).
    \end{equation*}
    Therefore, it remains to find the Jacobian relative to $x$ for our system:
    \begin{equation*}
        \frac{\delta f}{\delta x}([x_1, x_2]^T, u) = \mat{
            0 & 1 \\
            \frac{gm\ell}{I} \sin(x_1) & \frac{-b}{I}
        }
    \end{equation*}
    Thus our linearization (neglecting higher order terms) is given by:
    \begin{align*}
        \frac{d}{dt}(\delta x^*) &\approx A_1 \delta x^* \\
        \frac{d}{dt}(\delta x^{**}) &\approx A_2 \delta x^{**} \\
    \end{align*}
    where $A_1, A_2\in\R^{2\times 2}$ are
    \begin{align*}
        A_1 &= \frac{\delta f}{\delta x}(x^*, u^*) = \mat{
            0 & 1 \\
            \frac{gm\ell}{I} \sin(-\pi/2) & \frac{-b}{I}
        } = \mat{
            0 & 1 \\
            -\frac{gm\ell}{I} & \frac{-b}{I}
        } \\
        A_2 &= \frac{\delta f}{\delta x}(x^{**}, u^*) = \mat{
            0 & 1 \\
            \frac{gm\ell}{I} \sin(\pi/2) & \frac{-b}{I}
        } = \mat{
            0 & 1 \\
            \frac{gm\ell}{I} & \frac{-b}{I}
        }.
    \end{align*}

\item By Lyapunov's first method, we have that
\begin{enumerate}
    \item if the linearization is (globally) exp. stable,
        the nonlinear is locally exp. stable;
    \item if the linearization is (globally) unstable,
        the nonlinear is locally unstable; and
    \item if the linearization is (globally) marginally,
        the method is indecisive.
\end{enumerate}

So we analyze the real components of the eigenvalues of the matrices
$A_1$ and $A_2$ to determine the local stability properties of $x^*$ and
$x^{**}$, respectively.

Starting with the first equilibrium point $x^*$, and assuming
$\ell = m = I = 1$, $g = 10$, and $b \geq 0$, we have
\begin{equation*}
    A_1 = \mat{
        0 & 1 \\
        -10 & -b
    }
\end{equation*}
which has eigenvalues given by the characteristic equation:
\begin{align*}
    -\lambda (-b - \lambda) + 10 = 0 \\
    \lambda^2 + b\lambda + 10 = 0 \\
    \lambda_{1,2} = \frac{1}{2} (-b \pm \sqrt{b^2 - 40})
\end{align*}
where $\lambda_{1,2}$ are the first and second eigenvalues (corresponding
to the (+) and (-) of the $\pm$, respectively).

For $b = 0$, $\lambda_{1,2} = \pm i \sqrt{40}$ with $i = \sqrt{-1}$. So the 1st
Lyapunov method is indecisive.

For $0 < b^2 \leq 40$, $\Re(\sqrt{b^2 - 40}) = 0$, so that 
\begin{equation*}
    \Re(\lambda_{1,2}) = -\frac{b}{2} < 0
\end{equation*}

For $b^2 > 40$, the square root term is real and since $b^2 - 40 < b^2$,
we have $\sqrt{b^2 - 40} < |b| = b$ and thus
\begin{align*}
    \Re(\lambda_{1,2}) &= 1/2 \Re(-b \pm \sqrt{b^2 - 40}) \\
        &= 1/2 (-b \pm \sqrt{b^2 - 40}) \\
        &< 0
\end{align*}

So in summary, for any $b > 0$, we have $\Re(\lambda_{1,2}) < 0$ and thus the system is
locally exponentially stable about $x^*$. And again, for $b=0$ Lyapunov's first method
is indecisive.

Repeating the analysis for $x^{**}$ (and skipping steps which are nearly identical),
we have
\begin{equation*}
    A_2 = \mat{
        0 & 1 \\
        10 & -b
    }
\end{equation*}
with eigenvalues
\begin{equation*}
    \lambda_{1,2} = 1/2 (-b \pm \sqrt{b^2 + 40}).
\end{equation*}
Since $b^2 + 40 > 0$ for all $b$, the root is strictly real, and we have
$\sqrt{b^2 + 40} > b$ so that
\begin{equation*}
    \Re(\lambda_1) = \frac{1}{2} \Re(-b + \sqrt{b^2 + 40}) > 0.
\end{equation*}
The linearized system has a positive eigenvalue and is (globally)
unstable. Therefore, the nonlinear system is locally unstable about
$x^{**}$ for all values of $b$ (even negative values, but of course including
all those non-negative).
\end{letterenum}

\section{Problem 2}
Let $\dot{x} = Ax + Bu$ where
\begin{equation*}
    A = \mat{
        -2 & 1 \\
        0 & 0
    } \qquad B = \mat{
        0 \\ 1
    }
\end{equation*}
The solution to the non-homogenous LTI system is given by:
\begin{equation}
    x(t) = e^{At} x(0) + \int_{0}^{t} e^{A (t - \tau)} B u(\tau) d\tau
    \label{eq:non-homogenous-solution}
\end{equation}

To facilitate computation of the exponential $e^{At}$, we first compute
its jordan normal form. $A$'s characteristic equation is given by:
\begin{align*}
    (-2 - \lambda) (-\lambda) = \lambda( \lambda + 2) &= 0 \\
    \lambda_{1,2} &= 0, -2
\end{align*}
Thus, $A$ is diagonizable, and it is easy to see that corresponding eigenvectors
are $v_1 = [1, 2]^T$ and $v_2 = [1, 0]^T$. We can therefore write
$A$ in the form $A = PDP^{-1}$ where
\begin{equation*}
    P = \mat{1 & 1 \\ 2 & 0} \qquad D = \mat{0 & 0 \\ 0 & -2}
        \qquad P^{-1} = \frac{1}{2} \mat{0 & 1 \\ 2 & -1}
\end{equation*}
and it follows that $e^{At} = Pe^{Dt}P^{-1}$.

From equation (\ref{eq:non-homogenous-solution}), with initial state
$x(0) = 0_{2 \times 1}$ and input $u(t) = e^{-t}$, we have
\begin{align*}
    x(t) &= e^{At} \int_{0}^{t} e^{-A\tau} B e^{-\tau} d\tau \\
        &= P e^{Dt} P^{-1} \int_{0}^{t} P e^{-D\tau} (P^{-1} B) e^{-\tau} d\tau \\
        &= \frac{1}{2} P e^{Dt} \int_{0}^{t} \mat{
            1 & 0 \\
            0 & e^{2\tau}
        } \mat{1 \\ -1}e^{-\tau} d\tau \\
        &= \frac{1}{2} P e^{Dt} \int_{0}^{t} \mat{e^{-\tau} \\ -e^{\tau}} d\tau \\
        &= \frac{1}{2} P e^{Dt} \mat{-e^{-\tau} \\ -e^{\tau}} \eval_{\tau=0}^t \\
        &= \frac{1}{2} P \mat{
            1 & 0 \\
            0 & e^{-2t}
        } \biggl(
            \mat{-e^{-t} \\ -e^{t}} - \mat{-1 \\ -1}
        \biggr) \\
        &= \frac{1}{2} P \mat{
            1 - e^{-t} \\
            e^{-2t} ( 1- e^{t})
        } \\
        &= \frac{1 - e^{-t}}{2} \mat{1 & 1 \\ 2 & 0}\mat{1 \\ -e^{-t}} \\
        &= \frac{1-e^{-t}}{2} \mat{1 - e^{-t} \\ 2}
\end{align*}
Therefore, at $t=1$, we have
\begin{equation*}
    x(1) = \frac{1 - e^{-1}}{2} \mat{1 - e^{-1} \\ 2} \approx \mat{0.1998 \\ 0.6321}
\end{equation*}

\section{Problem 3}
We begin this problem by proving that a homogenous LTV system $\dot{x} = A(t) x$
has solution
\begin{align}
    x(t) = \Phi(t, t_0) x(t_0) \\
    \Phi(t, t_0) = e^{\int_{t_0}^{t} A(\tau) d\tau}
    \label{eq:ltv-solution}
\end{align}
if $A(t)$ and $\int_{t_0}^{t} A(\tau) d\tau$ commute for all $t$.

For shorthand throughout the exercise, we will define the function
$f_A: \R \to \R^n$ such that $f_A(t)$ is equal
to the integral term which appears in the exponent of equation (\ref{eq:ltv-solution}):
$f_A(t) = \int_{t_0}^{t} A(\tau) d\tau$.

We will prove that (\ref{eq:ltv-solution}) holds, given $A(t)$ and $f_A(t)$ commute.
First, we prove that given $A(t)$ and $f_A(t)$
commute, the following relation holds for any $i\in\Z_{\geq0}$:
\begin{equation}
    \frac{d}{dt} \biggl(
        \biggl(
            \int_0^t A(\tau) d\tau
        \biggr)^i
    \biggr) = i A(t) \biggl(\int_{0}^{t} A(\tau) d\tau\biggr)^{i-1}
    \label{eq:derivative-of-ltv-integral}
\end{equation}

\begin{proof}
    We will prove by induction. The base case $i=0$ clearly holds since
    $\frac{d}{dt}(1) = 0$.  Then, given that (\ref{eq:derivative-of-ltv-integral})
    holds for $i$, for $i+1$, we have
    \begin{align*}
        \frac{d}{dt} \biggl(
            f_A(t)^{i+1}
        \biggr)
        &= \frac{d}{dt} (f_A(t) f_A(t)^i) \\
        &= f_A(t) (i A(t) f_A(t)^{i-1}) + A(t) f_A(t)^{i} \\
        &= i A(t) f_A(t)^{i} + A(t) f_A(t) \\
        &= (i + 1) A(t) f_A(t)^{i}
    \end{align*}
    where we have used the product rule, the fact that $\frac{df_A}{dt}(t) = A(t)$,
    and the fact that $A(t)$ and $f_A(t)$ commute.
    Thus, (\ref{eq:derivative-of-ltv-integral})
    holds for all $i$.
\end{proof}

Now, since $x(0)$ is
constant, we are interested in the derivative of the exponential component.
Since matrix exponentiation is defined as $e^A := \sum_{i=0}^{\infty} A^i / i!$
for any $A \in \R^{n \times n}$, by applying (\ref{eq:derivative-of-ltv-integral})
we obtain the following
\begin{align*}
    \frac{d}{dt} (e^{f_A(t)}) &= \sum_{i=0}^{\infty} \frac{\frac{d}{dt} f_A(t)^i}{i!} \\
        &= \sum_{i=1}^{\infty} \frac{i A(t) f_A(t)^{i-1}}{i!} \\
        &= A(t) \sum_{i=1}^{\infty} \frac{f_A(t)^{i-1}}{(i-1)!} \\
        &= A(t) e^{f_A(t)}
\end{align*}

Finally, differentiating the solution proposed in (\ref{eq:ltv-solution}) yields:
\begin{align*}
    \frac{d}{dt}(x(t)) = \frac{d}{dt}(e^{\int_{0}^{t} A(\tau) d\tau})
        = A(t) e^{\int_{0}^{t} A(\tau)d\tau} x(0)
        = A(t) x(t)
\end{align*}

To see that the initial condition $x(t_0)$ holds, note that
$f_A(t_0) = 0_{n \times n}$ so that
$\Phi(t_0, t_0) = e^{0_{n\times n}} = I_{n\times n}$.

Returning to the case where
\begin{equation*}
    A(t) = \mat{0 & \omega(t) \\ -\omega(t) & 0},
\end{equation*}
we see that $A(t)$ and $A(\tau)$ commute by simply carrying out the matrix
multiplication and using the fact that the real numbers $\omega(t)$ and
$\omega(\tau)$ commute:
\begin{align*}
    A(t) A(\tau) &= \mat{
        0 & \omega(t) \\
        -\omega(t) & 0
    } \mat{
        0 & \omega(\tau) \\
        -\omega(\tau) & 0
    } \\
    &= \mat{
        -\omega(t) \omega(\tau) & 0 \\
        0 & -\omega(t) \omega(\tau)
    } \\
    &= \mat{
        0 & \omega(\tau) \\
        -\omega(\tau) & 0
    } \mat{
        0 & \omega(t) \\
        -\omega(t) & 0
    } \\
    &= A(\tau) A(t).
\end{align*}

Now, multiplication of $A(t)$ and $f_A(t)$ gives

\begin{align*}
    A(t) f_A(t) &= A(t) \int_{0}^{t} A(\tau) d\tau \\
        &= \int_{0}^{t} A(\tau) d\tau A(t) \\
        &= f_A(t) A(t)
\end{align*}
where we have pulled $A(t)$ into the integral and used the fact that
$A(t)$ and $A(\tau)$ commute.

Thus, the equation (\ref{eq:ltv-solution}) provides a valid solution, and
it remains to simplify the integral $f_A(t)$. Since $\omega(t)$ is a constant
we can factor it out of the matrix. Let $B = \mat{0 & 1 \\ -1 & 0}$ such
that $A = \omega(t)B$. Since the eigenvalues of $B$ are $\pm i$, we use
the following diagonalization $B = PDP^{-1}$ where
\begin{equation*}
    P = \mat{
        1 & 1 \\
        i & -i
    } \qquad D = \mat{
        i & 0 \\
        0 & -i
    } \qquad P^{-1} = \frac{1}{2} \mat{
        1 & -i \\
        1 & i
    }
\end{equation*}

Define $W(t) = \int_{0}^{t} \omega(\tau) d\tau$. Then,
\begin{align*}
    e^{f_A(t)} &= e^{\int_{0}^{t} \omega(\tau) PDP^{-1}d\tau} \\
        &= e^{P (W(t) D) P^{-1}} \\
        &= P e^{W(t) D} P^{-1} \\
        &= \mat{
            1 & 1 \\
            i & -i
        } \mat{
            e^{iW(t)} & 0 \\
            0  & e^{-iW(t)}
        } \frac{1}{2} \mat{
            1 & -i \\
            1 & i
        } \\
        &= \frac{1}{2} \mat{
            e^{iW(t)} & e^{-iW(t)} \\
            i e^{iW(t)} & -i e^{-iW(t)}
        } \mat{
            1 & -i \\
            1 & i
        } \\
        &= \frac{1}{2} \mat{
            e^{iW(t)} + e^{-iW(t)} & -i(e^{iW(t)} - e^{-iW(t)}) \\
            i(e^{iW(t)} - e^{-iW(t)}) & e^{iW(t)} + e^{-iW(t)}
        }.
\end{align*}
By Euler's formula, for any $a \in \R$ we have
\begin{equation*}
    e^{ia} + e^{-ia} = (\cos a + i \sin a) + (\cos a - i \sin a) = 2\cos a
\end{equation*}
and 
\begin{equation*}
    i(e^{ia} - e^{-ia}) = i(\cos a + i\sin a) - i(\cos a - i \sin a) = 2i^2\sin a = -2 \sin a
\end{equation*}
so that
\begin{equation}
    \Phi(t, t_0) = e^{f_A(t)} = \mat{
        \cos W(t) & \sin W(t) \\
        -\sin(W(t)) & \cos W(t)
    }.
    \label{eq:exponential-of-skew-integral}
\end{equation}

\section{Problem 4}

\begin{letterenum}

\item The derivative of $\frac{1}{2}\norm{x(t)}^2$ can be computed as follows
\begin{align*}
    \frac{d}{dt}\biggl(\frac{1}{2}\norm{x(t)}^2\biggr)
        &= \frac{d}{dt} \biggl(\frac{1}{2}x(t)^T x(t)\biggr) \\
        &= \frac{1}{2} \biggl(\dot{x}^T(t) x(t) + x^T(t) \dot{x}(t)\biggr) \\
        &= x^T(t) \dot{x}(t) \\
        &= x^T(t) A(t) x(t) \\
\end{align*}
where the third equality follows from the fact that the vector inner product
commutes ($\dot{x}^T(t) x(t) = x^T(t) \dot{x})$. Let $x(t) = [x_1, x_2, x_3]^T$,
where we have removed the dependence on time simply as a notational shorthand.

Substituting our given matrix $A(t)$ into the previous relation yields:
\begin{align*}
    \frac{d}{dt}\biggl(\frac{1}{2} \norm{x(t)}^2\biggr) &=
        \mat{x_1 & x_2 & x_3} \mat{
            t x_2 - 2x_3 \sin (2t) \\
            -t x_1 + x_3 e^t \\
            2x_1 \sin (2t) - x_3 e^t
        } \\
        &= t(x_1x_2 - x_1x_2) + 2\sin(2t) (x_1x_3 - x_1x_3) + e^t (x_2x_3 - x_2x_3) \\
        &= 0
\end{align*}

Therefore, since $d/dt(\frac{1}{2}\norm{x(t)}^2) = 0$ for all $t$,
integrating both sides yields
\begin{equation*}
    \norm{x(t)}^2 = C
\end{equation*}
where $C$ is a integration constant, easily seen to be $C = \norm{x(0)}^2$.
Since $\norm{x} \geq 0$ for any $x\in\R^n$ we have:
\begin{equation}
    \norm{x(t)} = \norm{x(0)}
    \label{eq:constant-norm}
\end{equation}

\item We first note that since the system is linear, it has one
equilibrium point located at the origin, $x^* = 0$.
Then, for any $\delta > 0$, if we select $x(0)$ such that $\norm{x(0)} < \delta$,
by equation (\ref{eq:constant-norm}) we have $\norm{x(t)} = \norm{x(0)} < \delta$
for all $t$. Selecting $\epsilon = \delta$, we therefore have that the norm
of $x(t)$ is bounded by $\epsilon$ for all $t$, and $x^*$ satisfies the definition
of a marginal stability equilibrium point.

However, since $\lim_{t\to\infty} \norm{x(t)} = \norm{x(0)}$, for any
$x(0)$ with $\norm{x(0)} > 0$
\begin{equation*}
    \lim_{t\to\infty} \norm{x(t)} \neq 0
\end{equation*}
implies that 
\begin{equation*}
    \lim_{t\to\infty} x(t) \neq \mat{0 \\ 0 \\ 0} = x^*
\end{equation*}
by the definition of a norm ($\norm{x} = 0$ iff $x = 0$).

Since $x(t)$ does not converge to the origin, the origin is not a
asymptotically (or exponentially) stable equilibrium point.





\end{letterenum}


\end{document}
